\documentclass[11pt]{article}
\usepackage{fontspec}
\usepackage[french]{babel}
\usepackage{amsmath}
\usepackage{amssymb}
\usepackage{amsthm}
\usepackage{mathrsfs}
\usepackage{tikz-cd}
\usepackage{enumitem}

\newcommand{\printedpage}[1]{\marginpar{\footnotesize p.~#1}}
\newcommand{\msp}[1]{\marginpar{\footnotesize [#1]}}
\DeclareMathOperator{\card}{card}
\DeclareMathOperator{\Ob}{Ob}
\DeclareMathOperator{\Ind}{Ind}

\begin{document}

\printedpage{124}
% head: A. Grothendieck : Dérivateurs, Ch. XVIII, Édition des Universités de Montpellier II et Paris VII

compatibilités $3_\alpha(\beta)$, donc des double-flèches
\[
\begin{tikzcd}
\varphi(i_\beta) \arrow[r, shift left, "\tau'_{\alpha,\beta}u_\beta"] \arrow[r, shift right, "u_\alpha\varphi(\tau_{\alpha,\beta})"'] & j_\alpha \; .
\end{tikzcd}
\]
Comme l'ensemble de ces double-flèches est de cardinal $\card\alpha < \pi$, et que $J'$ est grande devant ledit cardinal, on
\msp{page 329}
trouve donc une flèche
\[
\rho : j_\alpha \longrightarrow j
\]
dans $J$ qui les égalise toutes. Comme $J'$ est cofinal, on peut supposer $j$ dans $J'$, soit $\bar{j}_\alpha$. Donc quitte à remplacer $j_\alpha$ par $\bar{j}_\alpha$, et les $\tau'_{\alpha,\beta}$ par $\bar{\tau}'_{\alpha,\beta} = \rho\tau'_{\alpha,\beta}$, enfin $u_\alpha : \varphi(i_\alpha) \longrightarrow j_\alpha$ par $\bar{u}_\alpha = \rho u_\alpha : \varphi(i_\alpha) \longrightarrow \bar{j}_\alpha$, on satisfait à toutes les contraintes voulues.

Je résume le résultat (``$\pm$ tautologique''!) prouvé dans les pages 313-329.

% a vertical change bar is printed in the left margin alongside the statement from "Lemme-clef" to "suivantes :"
\textbf{Lemme-clef 4.13.17.} \textit{Considérons un carré cartésien dans} Cat
\[
\begin{tikzcd}
I \arrow[r, "\varphi"] & J \\
I' \arrow[u, hook] \arrow[r] & J' \arrow[u, hook, "\psi"'] \; .
\end{tikzcd}
\]

\textit{On fait les hypothèses suivantes :}

\msp{page 330}

\begin{enumerate}[label={}]
\item[\textit{1)}] \textit{$\varphi$ et $\psi$ sont cofinaux, $\psi$ est l'inclusion d'une sous-catégorie strictement pleine $J'$ de $J$, $I$ est filtrant (donc $J$ et $J'$ aussi).}
\end{enumerate}

\textit{On cherche des conditions pour que l'inclusion de la sous-catégorie strictement pleine $I' \longrightarrow I$ soit cofinale (auquel cas on voit tout de suite que $I'$ est filtrante, et $I' \longrightarrow J'$ est aussi cofinal). Pour ceci, on suppose qu'il existe un \textup{ordinal limite} $\alpha_0$ ayant les propriétés supplémentaires suivantes :}

\begin{enumerate}[label={}]
\item[\textit{2)}] \textit{$I$ et $J'$ sont stables par $\varinjlim$ de type $I_{\alpha_0}$ $(= \{\beta \mid \beta \ \textit{ordinal}\ < \alpha_0\})$, et les foncteurs $\varphi$, $\psi$ commutent à ces limites.}
\item[\textit{3)}] \textit{Pour tout $\alpha \in I_{\alpha_0}$, $I$ est grande devant $\card\alpha$. (Donc, $\varphi$ étant cofinal, $J$ est également grande devant $\card\alpha$, et de même, $\psi$ étant pleinement fidèle, et cofinal, $J'$ est grande devant $\card\alpha$.)}
\end{enumerate}

\textit{Sous ces conditions, $I'$ est cofinal dans $I$ (donc aussi grande devant $\card\alpha$, pour tout $\alpha \in I_{\alpha_0}$).}

La démonstration est celle donnée ci-dessus. On n'a pas eu à utiliser que l'ordinal $\alpha_0 = \alpha_\pi$ était de cette espèce particulière, seulement que c'était un cardinal limite.

\printedpage{125}
% head: A. Grothendieck : Dérivateurs, Ch. XVIII, Édition des Universités de Montpellier II et Paris VII

\msp{page 331}

Ce lemme implique, comme on l'a vu à la page 312, l'implication b $\Longrightarrow$ d dans le diagramme d'implications $(*)$ page 310, lequel est donc à présent prouvé.

Pour achever de prouver le théorème 4.13.15, il reste à prouver les implications
\[
\begin{tikzcd}[row sep=small]
\mathrm{c}' \arrow[dr, Rightarrow] & & \\
& \mathrm{a}' \arrow[r, Rightarrow] & \mathrm{a} \; . \\
\mathrm{b} \arrow[ur, Rightarrow] & &
\end{tikzcd}
\]

Démontrons donc (b ou c$'$) $\Longrightarrow$ a$'$. Il reste à prouver que $\underline{S}$ est $\pi$-accessible dans $\mathcal{M}$, i.e. que le foncteur pleinement fidèle
\[
\Ind_\pi(\underline{S}(\pi)) \longrightarrow \mathcal{M} \tag*{$(*)$}
\]
a pour image essentielle $\underline{S}$. Donc il faut prouver que si $x \in \Ob\mathcal{M}$, alors
\[
x \in \Ob\underline{S} \quad \overset{(?)}{\Longleftrightarrow} \quad
\begin{minipage}[c]{0.55\textwidth}
$x$ est dans l'image essentielle de $\Ind_\pi(\underline{S}(\pi))$, i.e. $x$ est isomorphe à $\varinjlim\limits_{i\in I} x_i$, les $x_i$ dans $\Ob\underline{S}(\pi)$, $I$ grande devant $\pi$.
\end{minipage}
\tag*{$(*)$}
\]

\msp{page 332}

$\Bigg($ Supposons d'abord que $\underline{S}$ satisfasse b) (\footnote[141]{Maladroit, on peut traiter les deux cas b) et c) en même temps, cf. bas de page 333.}). Donc on sait que
\[
x \in \Ob\underline{S} \qquad \Longleftrightarrow \qquad \underline{S}/x \text{ est cofinale dans } \mathcal{M}_\pi/x \; . \tag*{$(**)$}
\]

Nous voulons déduire l'équivalence $(*)$ de l'équivalence $(**)$.

Prouvons $\Longrightarrow$ dans $(*)$, i.e. supposons $x \in \Ob\underline{S}$, et prouvons que $x$ est dans l'image essentielle de $\Ind_\pi\underline{S}(\pi)$. En vertu de $(**)$, $\underline{S}(\pi)/x$ est une sous-catégorie (strictement pleine) cofinale dans $\mathcal{M}_\pi/x$, donc comme $\mathcal{M}_\pi/x$ est grande devant $\pi$, $\underline{S}(\pi)/x$ l'est aussi. Comme on sait que
\[
x \ = \ \varinjlim_{i\in\mathcal{M}_\pi/x} x_i
\]
($\pi$ étant cardinal utile pour $\mathcal{M}$), il s'ensuit par cofinalité qu'on a aussi
\[
x \ = \ \varinjlim_{i\in\underline{S}(\pi)/x} x_i \; ,
\]
et comme $\underline{S}(\pi)/x$ est grande devant $\pi$, cela implique que $x$ appartient à l'image essentielle de $\Ind_\pi(\underline{S}(\pi))$.

Prouvons $\Longleftarrow$ dans $(*)$, i.e. supposons que $x$ soit dans l'image essentielle de $\Ind_\pi(\underline{S}(\pi))$,
\printedpage{126}
% head: A. Grothendieck : Dérivateurs, Ch. XVIII, Édition des Universités de Montpellier II et Paris VII
\msp{page 333}
\[
x \ = \ \varinjlim_{i\in I} x_i \; , \qquad\quad \text{les } x_i \text{ dans } \underline{S}(\pi),\ I \text{ grande devant } \pi.
\]

Alors le foncteur composé $i \longmapsto x_i$,
\[
I \ \longrightarrow \ \underline{S}(\pi)/x \ \hookrightarrow \ \mathcal{M}_\pi/x
\]
est cofinal (cf. 4.4.2, page 79), donc $\underline{S}(\pi)/x \longrightarrow \mathcal{M}_\pi/x$ satisfait la condition F 1 (puisque le composé le satisfait), et comme $\underline{S}(\pi)/x \longrightarrow \mathcal{M}_\pi/x$ est pleinement fidèle, cela implique déjà que ce foncteur est cofinal. Donc par $(**)$ on en conclut $x \in \Ob\underline{S}$, q.e.d.

Supposons maintenant que $S$ satisfasse c$'$ (au lieu de b), et prouvons encore l'équivalence $(*)$. Pour ceci, écrivons
\[
\square \qquad x \ = \ \varinjlim_{i\in I} x_i \qquad\quad
\left\{
\begin{minipage}[c]{0.5\textwidth}
les $x_i \in \Ob\mathcal{M}_\pi$, $I$ \textit{ensemble ordonné}, $(x_i)_{i\in I}$ est $\pi_0$-$\pi$-adaptée (\footnotemark[142]).
\end{minipage}
\right.
\]
\footnotetext[142]{\textbf{N.B.} Ça existe, grâce à 4.13.14.}

ce qui est possible puisque $\pi$ est utile pour $\mathcal{M}$ $\Bigg)$. Définissons $I_{\underline{S}}$ comme dans l'énoncé de c), donc comme image inverse de $\underline{S}(\pi)/x \hookrightarrow \mathcal{M}_\pi/x$ par $I \longrightarrow \mathcal{M}_\pi$. Par hypothèse d$'$) \texttt{[plutôt c$'$)]}, on a
\[
x \in \Ob\underline{S} \qquad \Longleftrightarrow \qquad I_{\underline{S}} \text{ est cofinal dans } I \; . \tag*{$(**')$}
\]

En utilisant $(**')$, prouvons encore $(*)$ (page 331). En fait, on utilisera seulement qu'on a
\msp{page 334}
l'isomorphisme $\square$, avec $I$ grande devant $\pi$, et qu'on a l'équivalence $(**')$, sans avoir à utiliser que $I$ est ordonnée. Donc la démonstration vaut aussi dans le cas c).

Prouvons $\Longrightarrow$ dans $(*)$, i.e. supposons $x \in \Ob\underline{S}$, prouvons que $x$ est dans l'image essentielle de $\Ind_\pi(\underline{S}(\pi))$. En vertu de $(**')$, $I_{\underline{S}}$ est cofinale dans $I$, et comme $I_{\underline{S}} \longrightarrow I$ est une inclusion pleinement fidèle, il s'ensuit que $I_{\underline{S}}$ est également grande devant $\pi$, d'autre part on aura par cofinalité
\[
x \ = \ \varinjlim_{i\in I_{\underline{S}}} x_i \; ,
\]
où cette fois les $x_i$ sont dans $\underline{S}(\pi) = \underline{S}\cap\mathcal{M}_\pi$, d'où résulte que $x$ est dans l'image essentielle de $\Ind_\pi(\underline{S}(\pi))$.

Prouvons $\Longleftarrow$ dans $(*)$, i.e. supposons que l'on ait
\[
x \ = \ \varinjlim_{j\in J} x_j \; , \qquad\quad \text{les } x_j \text{ dans } \underline{S}(\pi),\ J \text{ grande devant } \pi \; .
\]

\printedpage{127}
% head: A. Grothendieck : Dérivateurs, Ch. XVIII, Édition des Universités de Montpellier II et Paris VII

\msp{page 335}

En invoquant à nouveau 4.4.2, on trouve que $J \longrightarrow \mathcal{M}_\pi/x$ est cofinal, et comme ce foncteur se factorise par le foncteur d'inclusion pleinement fidèle $\underline{S}(\pi)/x \hookrightarrow \mathcal{M}_\pi/x$, il s'ensuit comme tantôt que ce dernier est cofinal. Nous allons admettre pour un moment que ceci implique que l'image inverse $I_{\underline{S}}$ de cette sous-catégorie strictement pleine, par le foncteur également cofinal $I \longrightarrow \mathcal{M}_\pi/x$, est cofinal dans $I$. Ceci étant, l'équivalence $(**')$ (p. 333) implique $x \in \Ob\underline{S}$, ce qu'on voulait.

Donc il reste à prouver que si
\[
J' \ = \ \underline{S}(\pi)/x \hookrightarrow \mathcal{M}_\pi/x \ = \ J
\]
est cofinal, alors $I' = I_{\underline{S}} \longrightarrow I$ l'est aussi. Il semble qu'il faille bel et bien à nouveau invoquer les conditions de stabilité $\mathrm{L}(\pi_0,\pi)$ faites dans d$'$) \texttt{[plutôt c$'$)]} ce qui n'avait pas été nécessaire pour le cas c) et
\msp{page 336}
le lemme-clef (un peu) vache 4.13.17, où on prendra (comme dans la démonstration écrite
avant l'énoncé de cet lemme) % [sic] cet lemme -> ce lemme, demonstrative adjective "cet" before a consonant
$\alpha = \alpha_\pi \overset{\text{déf}}{=}$ plus petit ordinal $\alpha$ tel que $\card\alpha = \pi$. Donc $I_\alpha$ est de cardinal $\pi$, et grande devant tout cardinal $\pi' < \pi$, et en particulier devant $\pi'_0$. L'hypothèse $2^\circ$) de 4.13.17 est vérifiée pour $I$, par définition d'un système inductif $\pi_0$-$\pi$-adapté, et pour $J'$, grâce à l'hypothèse $\mathrm{L}(\pi_0,\pi)$ sur $\underline{S}(\pi)$. De même, l'hypothèse $3^\circ$) ($I$ grande devant tout cardinal $< \pi$) est satisfaite par définition des systèmes $\pi_0$-$\pi$-adaptés -- $I$ est même grande devant $\pi$. L'hypothèse $1^\circ$) a déjà été vue.

On a donc prouvé (par voie circulaire)
\[
\begin{tikzcd}[row sep=small]
& & & \mathrm{c}' & \\
\mathrm{a} \arrow[r, Rightarrow] & \Bigg(\mathrm{a}' \arrow[r, Leftrightarrow] & \mathrm{c} \arrow[ur, Leftrightarrow] \arrow[dr, Rightarrow] & & \Bigg)\, , \\ % Verified against PDF leaf 9: the oblique equivalence has both arrowheads.
& & & \mathrm{b} &
\end{tikzcd}
\]

\texttt{[et aussi b $\Longrightarrow$ c,]} en ayant à invoquer le lemme-clef vache pour les \textit{deux} implications
\[
\mathrm{a}' \ \Longrightarrow \ \mathrm{c} \qquad\quad \text{et} \qquad\quad \mathrm{c}' \ \Longrightarrow \ \mathrm{a}'
\]
(donc pour deux implications \textit{en sens inverses}, chose étrange \dots).

\msp{page 337}

Il reste seulement encore à prouver l'implication
\[
\mathrm{a}' \ \Longrightarrow \ \mathrm{a} \; ,
\]
ce qui est essentiellement sorital.

Cela achève la démonstration du critère d'accessibilité 4.13.15 d'une sous-catégorie pleine $\underline{S}$ de $\mathcal{M}$.

On va donner une variante un peu décantée de 4.13.15, qui (avec des hypothèses un peu plus faibles) capte l'essentiel de sa substance. L'énoncé 4.13.15 en est une conséquence quasi immédiate :

\printedpage{128}
% head: A. Grothendieck : Dérivateurs, Ch. XVIII, Édition des Universités de Montpellier II et Paris VII

\msp{page 338}

\textbf{Corollaire 4.13.18} (de 4.13.15). \textit{Soient $\mathcal{M}$ une catégorie accessible, $\pi$ un cardinal utile pour $\mathcal{M}$, $\alpha_0$ un \textup{ordinal limite} tel que l'on ait $\card\alpha_0 \leqslant \pi$. Soit $I_{\alpha_0} = \{\alpha \mid \alpha < \alpha_0\}$, et supposons que $\mathcal{M}$ soit stable par $\varinjlim$ de type $I_{\alpha_0}$.}

\textit{Soit alors $\underline{S}$ une sous-catégorie strictement pleine de $\mathcal{M}$, $\underline{S}(\pi) = \underline{S}\cap\mathcal{M}_\pi$. Désignons par $L(\alpha_0)$ la condition : $\underline{S}(\pi)$ est stable dans $\mathcal{M}_\pi$ (ou ce qui revient au même, dans $\mathcal{M}$) par $\varinjlim$ de type $I_{\alpha_0}$. Alors les conditions suivantes sont équivalentes.}

\begin{enumerate}[label={}]
\item[\textit{a)}] \textit{$\underline{S}$ est $\pi$-accessible dans $\mathcal{M}$, et est stable dans $\mathcal{M}$ par $\varinjlim$ de type $I_{\alpha_0}$.}
\item[\textit{a$'$)}] \textit{$\underline{S}$ est $\pi$-accessible et $\underline{S}(\pi)$ satisfait $L(\alpha_0)$.}
\item[\textit{b)}] \textit{$\underline{S}(\pi)$ satisfait $L(\alpha_0)$, et pour tout $x \in \Ob\mathcal{M}$, on a ceci : $x \in \Ob\underline{S}$ si et seulement si la sous-catégorie strictement pleine $\underline{S}(\pi)/x$ de $\mathcal{M}(\pi)/x$ est cofinale dans $\mathcal{M}(\pi)/x$} \texttt{[lire $\mathcal{M}_\pi/x$]}.
\item[\textit{c)}] \textit{$\underline{S}(\pi)$ satisfait $L(\alpha_0)$, de plus pour tout $x \in \Ob\mathcal{M}$, et toute représentation}
\[
x \ \simeq \ \varinjlim_{i\in I} x_i
\]
\msp{page 339}
\textit{qui est ``$\alpha_0$-$\pi$-adaptée'', i.e. telle que}
\begin{enumerate}[label={}]
\item[\textit{$1^\circ$)}] \textit{$I$ équivalente à une petite catégorie, et grande devant $\pi$,}
\item[\textit{$2^\circ$)}] \textit{les $\varinjlim$ de type $I_{\alpha_0}$ existent dans $I$,}
\item[\textit{$3^\circ$)}] \textit{le foncteur $I \longrightarrow \mathcal{M}$, $i \longmapsto x_i$, commute aux $\varinjlim$ de type $I_{\alpha_0}$, et enfin}
\item[\textit{$4^\circ$)}] \textit{les $x_i \in \Ob\mathcal{M}_\pi$}
\end{enumerate}
\textit{(comparer 4.13.11), on a ceci (où on désigne par $I_{\underline{S}}$ la sous-catégorie strictement pleine de $I$ formée des $i$ tels que $x_i \in \Ob\underline{S}$) : on a $x \in \Ob\underline{S}$ si et seulement si $I_{\underline{S}}$ est cofinale dans $I$.}
\item[\textit{c$'$)}] \textit{(\textup{Si} $I_{\alpha_0}$ est grande devant $\pi_0$, et $\pi^{\pi_0} = \pi$, pour pouvoir invoquer 4.13.14, cf. p. 335.) Comme c), mais en supposant $I$ ordonnée.}
\end{enumerate}
(\footnote[143]{\textbf{N.B.} Si on ne suppose par $I_{\alpha_0}$ grande devant $\pi_0$ et % [sic] suppose par -> suppose pas, negation "ne ... pas"
$\pi^{\pi_0} = \pi$, alors toutes les conditions de a) à c) (a, a$'$, b, c) sont équivalentes et \textit{impliquent} c, mais ne lui sont peut-être pas équivalentes.}).

\textsc{Démonstration}. Elle se fait suivant le même diagramme logique que celle de 4.13.15,
\[
\begin{tikzcd}[row sep=small]
& & & \mathrm{c}' \arrow[dr, Rightarrow, "\textbf{!}"] & \nearrow\ \begin{tabular}[b]{@{}r@{}}\footnotesize ici, il faut supposer\\ \footnotesize $I_{\alpha_0}$ grande devant $\pi_0$\\ \footnotesize et $\pi^{\pi_0} = \pi$.\end{tabular} & \\
\mathrm{a} \arrow[r, Rightarrow] & \mathrm{a}' \arrow[r, Rightarrow, "\textbf{!}"] & \mathrm{c} \arrow[ur, Rightarrow] \arrow[dr, Rightarrow] & & \mathrm{a}' \arrow[r, Rightarrow] & \mathrm{a} \; , \\
& & & \mathrm{b} \arrow[ur, Rightarrow] & &
\end{tikzcd}
\]
où on a marqué du signe\,! les deux implications qui sont ``vaches'', en ce sens qu'elles font appel au lemme-clef 4.13.17 (pages 329-330), qui se prouve par une construction transfinie un peu vissée. La démonstration des sept implications du diagramme est essentiellement
celle donné dans les pages précédentes 311-337, % [sic] celle donné -> celle donnée, agreement with feminine "celle"

\end{document}
